\documentclass[10pt,a4paper]{amsart}
\usepackage[margin=19mm]{geometry}
\usepackage{amsmath,amssymb,mathtools}
\usepackage[scr=rsfs,cal=euler]{mathalfa}
\usepackage{xfrac}
\usepackage[unicode,hidelinks,bookmarks=false]{hyperref}
\newcommand{\ie}{i.e.\ }
\newcommand{\cf}{cf.\ }
\newcommand{\resp}{respectively\ }
\newcommand{\Subsection}{\subsection}
\providecommand{\nobreakdash}{\nobreak\hbox{-}}
\setlength{\emergencystretch}{2em}
\multlinegap0pt
\usepackage{amssymb}
\usepackage{mathtools}
\usepackage{enumerate}
\usepackage[shortlabels]{enumitem}
\usepackage{xcolor}
\newtheorem{cor}{Corollary}
\newtheorem{rem}{Remark}
\newtheorem{lemma}{Lemma}
\newtheorem{thm}{Theorem}
\newcommand{\N}{\mathbb{N}}
\newcommand{\R}{\mathbb{R}}
\newcommand{\T}{\mathbb{T}}
\def\heat#1{e^{#1 \Delta}}
\newcommand{\na}{\nabla}
\newcommand{\norm}[1]{\|#1\|}
\newcommand{\w}{\omega}
\title{Global well-posedness of the 3D Navier-Stokes\\ equations with helical $L^1$ vorticity}
\author{Francisco Gancedo, Antonio Hidalgo-Torn\'e}
\date{Source: arXiv:2609.05193v1 (CC BY 4.0)}
\begin{document}
\maketitle

\noindent\emph{Source and licence.} Global well-posedness of the 3D Navier-Stokes\\ equations with helical $L^1$ vorticity. Version arXiv:2609.05193v1 (CC BY 4.0), distributed by arXiv under the Creative Commons Attribution 4.0 International licence, http://creativecommons.org/licenses/by/4.0/, as recorded in the arXiv licence metadata for this exact version and shown on the abstract page https://arxiv.org/abs/2609.05193v1. Full licence notice: \texttt{licenses/arxiv-2609.05193-CC-BY-4.0.txt}.\par\smallskip

\noindent\textit{Editorial scope.} Counted unit: the article's thm thm:higherregrate (source0.tex lines 631--641). The article's complete original proof of it is delivered with it (source0.tex lines 643--721, one \texttt{proof} environment of 79 lines). No mathematics is written, repaired or re-derived: the statement and the proof are the article's own lines, in the article's own notation, and no retained line is substituted or re-flowed.  Retained uncounted as the source-local context the proof needs: lines 292--295; lines 390--399; lines 403--408; lines 415--418; lines 460--465; lines 482--484; lines 521--523; lines 888--910.  Nothing was omitted from any retained span, so the package carries no omission record.  No retained line contains a citation, so the wrapper carries no bibliography and no bibliography entry.  No retained line contains a figure environment, an image-inclusion command or drawing code, so no image is delivered and the package declares no asset dependency.  Editorial formatting only, in addition: an AMS \texttt{amsart} preamble that restates the article's own declarations and macros used by the retained lines, namely the environments \texttt{cor}, \texttt{rem}, \texttt{lemma}, \texttt{thm} on separate counters, together with the standard packages those lines need.  The retained environments print in this document as Corollary 1, Remark 1, Lemma 1, Theorem 1 in order of appearance, whereas the article prints the same items with the numbering of its own full document.  The complete native source of the article as distributed in the arXiv e-print for this exact version is bundled unchanged as \texttt{sources/209420/source0.tex}.
\begin{cor}\label{cor:intbetatruncated}
If $a,b\in \R$ and $-1<a$, then 
$$\int_{\frac{t}{3}}^t (t-\tau)^a \tau^b d\tau=t^{a+b+1}\int_{\frac{1}{3}}^1 (1-z)^a z^b dz\lesssim  t^{a+b+1}.$$
\end{cor}

\begin{rem}\label{uregularity} In particular, if $\omega$ has helical symmetry and $u=\nabla\wedge(-\Delta)^{-1}\omega$,
then, for $1<p<2$ and $\frac1q=\frac1p-\frac12$,
\[
\|u\|_{L^q}
 \lesssim \|\nabla u\|_{L^p}
 \lesssim \|\omega\|_{L^p},
\]
where the last inequality follows from the boundedness of the Riesz
transforms.
\end{rem}

\begin{lemma}\label{lem:gagliardoniremberg}
Let $f$ be a smooth vector field with helical symmetry. Then,
$$\|f\|^2_{L^\infty}\lesssim \|f\|_{L^4}\|\na f\|_{L^4},$$
and
$$\|f\|^2_{L^4}\lesssim \|f\|_{L^2}\|\na f\|_{L^2}.$$
\end{lemma}

\begin{lemma}\label{lemma:heathelicalestimates}
If $f$ has helical symmetry, $s\geq 0$, and $1\leq q\leq p< \infty$, then 
$$\|\Lambda^s e^{t\Delta}f\|_{L^p}\lesssim t^{\frac{1}{p}-\frac{1}{q}-\frac{s}{2}}\|f\|_{L^q}.$$
\end{lemma}

\begin{rem}\label{rem:heatsmallXphigherorder} By taking derivatives of the heat kernel we can show that, if $\w_0$ is an integrable function with helical symmetry,
\begin{equation}\label{eq:heatscaling}
t^{1-\frac{1}{p}+\frac{s}{2}}\|\Lambda^s\heat{t}\w_0\|_{L^p}\to 0 \quad \text{as}\quad t\to 0^+,   
\end{equation}
for any $s\geq 0$ and $1\leq p<\infty$, except for the case $(s,p)=(0,1).$
\end{rem}

\begin{equation}\label{eq:NSDuhamel}
\w(t)=\heat{t}\w_0+\int_0^t\na\heat{(t-\tau)}(\w\otimes u-u\otimes \w)(\tau)d\tau, \quad u=\na\wedge(-\Delta)^{-1}\w.
\end{equation}

\begin{thm}\label{thm:shorttimeWP}
Let $\w_0\in L^1$ be helically symmetric and divergence free. Then, there exists $T>0, C>0$ such that there is a unique mild solution $\omega$ to the Navier-Stokes equations \eqref{eq:NSDuhamel} in $C([0,T];L^1)\cap X^{4/3}_T $, with $\|\w\|_{X^{4/3}_T}<C$ and initial datum $\w_0$. Besides, $\|\w\|_{X^{4/3}_t}\to 0$ as $t\to 0^+$.
\end{thm}

\begin{lemma}\label{lem:KatoPonce}
Let $n_1,n_2\in \N\setminus{\{0\}}$, $\frac{1}{2}<p<\infty$, and
$$
1<q_1,q_2,r_1,r_2<\infty
$$
satisfying
$$
\frac{1}{p}
=
\frac{1}{q_1}+\frac{1}{r_1}
=
\frac{1}{q_2}+\frac{1}{r_2}.
$$
Given $s>\max(0, \frac{n_1+n_2}{p}-(n_1+n_2))$ or $s\in 2\mathbb{N}$, there exists a constant such that
for any $f,g\in \mathscr{S}(\R^{n_1}\times\T^{n_2})$ it holds
$$
\|\Lambda^s(fg)\|_{L^p}
\lesssim
\|\Lambda^s f\|_{L^{q_1}}\|g\|_{L^{r_1}}
+
\|f\|_{L^{q_2}}\|\Lambda^s g\|_{L^{r_2}}.
$$
\end{lemma}

\begin{thm}\label{thm:higherregrate}
Let $\w$ be the unique mild solution obtained in Theorem \ref{thm:shorttimeWP}, and $T$ its time of existence. Then, for any $s\geq 0$ and $1\leq p\leq  \infty$,
\begin{equation}\label{eq:scalingsolutionhigherregularity}
t^{1-\frac{1}{p}+\frac{s}{2}}\|\Lambda^s\w(t)\|_{L^p}\lesssim 1  
\end{equation}
for all $0<t\leq T.$ Besides, 
\begin{equation}\label{convergecetozero}
t^{1-\frac{1}{p}+\frac{s}{2}}\| \Lambda^s\w(t)\|_{L^p}\to 0 \text{ as }t\to 0, 
\end{equation}
except for the case $(s,p)=(0,1)$.
\end{thm}

\begin{proof}
The interpolation inequality
$$
t^{1-\frac{1}{p}}\|\omega(t)\|_{L^p}
\leq
\|\omega(t)\|_{L^1}^{\frac{4}{p}-3}
(t^{\frac14}\|\omega(t)\|_{L^{4/3}})^{4(1-\frac{1}{p})},
\qquad 1\leq p\leq \frac{4}{3},
$$
together with Theorem \ref{thm:shorttimeWP}, allows to extend the result to the case $s=0$ and $1<p<\frac{4}{3}$, thereby obtaining property \eqref{convergecetozero} as the same holds for $\norm{\w}_{X_t^{4/3}}$.

Next, we fix $4/3<p<2$, $0<t\leq T$, and $t/3<\sigma<2t/3,$ so that $\sigma\approx (t-\sigma)\approx t$. We take the $L^p$ norm in \eqref{eq:NSDuhamel}, and then we apply Lemma \ref{lemma:heathelicalestimates}, Theorem \ref{thm:shorttimeWP} and Remark \ref{uregularity} to deduce
\begin{equation}\label{eq:L2regularityestimate}
\begin{aligned}
\|\w(t)\|_{L^p}\leq& \|\heat{(t-\sigma)}\w(\sigma)\|_{L^p}+\int_\sigma^t\|\na\heat{(t-\tau)}(\w\otimes u-u\otimes \w)(\tau)\|_{L^p}d\tau\\
\lesssim & (t-\sigma)^{\frac{1}{p}-\frac{3}{4}}\|\omega(\sigma)\|_{L^{\frac{4}{3}}}+\int_\sigma^t (t-\tau)^{\frac{1}{p}-{1}-\frac{1}{2}}\|\w(\tau)\|_{L^\frac{4}{3}}\|u(\tau)\|_{L^4}d\tau \\
\lesssim & (t-\sigma)^{\frac{1}{p}-\frac{3}{4}}\sigma^{-\frac{1}{4}}\|\omega\|_{X_t^{\frac{4}{3}}}+\|\omega\|^2_{X_t^{\frac{4}{3}}}\int_\sigma^t (t-\tau)^{\frac{1}{p}-\frac{3}{2}}\tau^{-\frac{1}{4}}\tau^{-\frac{1}{4}}d\tau\\
\lesssim & t^{\frac{1}{p}-1}\|\omega\|_{X_t^{\frac{4}{3}}}+ t^{\frac{1}{p}-1}\|\omega\|^2_{X_t^{\frac{4}{3}}}\int_{\frac{1}{3}}^1 (1-z)^{\frac{1}{p}-\frac{3}{2}}z^{-\frac{1}{2}}dz.
\end{aligned}
\end{equation}
The inequality above provides the bound of $\norm{\w}_{X_t^p}$ for any fixed $4/3<p<2$, and this norm goes to 0 as $t\to0^+$.

We now extend the estimate to any $2<p<\infty$.
The linear part can be estimated as in \eqref{eq:L2regularityestimate}. For the nonlinear term, to ensure the boundedness of the time integral, we choose $q=\frac{2p}{p+1}$ when applying Lemma \ref{lemma:heathelicalestimates}. This gives
$$
\begin{aligned}
\|\w(t)\|_{L^p}\leq& \|\heat{(t-\sigma)}\w(\sigma)\|_{L^p}+\int_\sigma^t\|\na\heat{(t-\tau)}(\w\otimes u-u\otimes \w)(\tau)\|_{L^p}d\tau\\
\lesssim & t^{\frac{1}{p}-1}\|\omega\|_{X_t^{\frac{4}{3}}}+\int_\sigma^t (t-\tau)^{\frac{1}{p}-\frac{p+1}{2p}-\frac{1}{2}}\|\w(\tau)\|_{L^\frac{4p}{2p+1}}\|u(\tau)\|_{L^{4p}}d\tau\\
\lesssim & t^{\frac{1}{p}-1}\|\omega\|_{X_t^{\frac{4}{3}}}+ t^{\frac{1}{p}-1}\|\w\|^2_{X_t^\frac{4p}{2p+1}}\int_{\frac{1}{3}}^1 (1-z)^{\frac1{2p}-1}z^{\frac{1}{2p}-1}dz.
\end{aligned}$$
As before, we obtain the desired bound.
By interpolation, the result for $p=2$ is also obtained, concluding the proof for $s=0$ and $1\leq p<\infty.$

We now turn to higher-order derivatives. Lacking a priori estimates for derivative norms, we cannot directly include the target norm on the right-hand side of our estimates. Crucially, applying full derivatives to the nonlinear term yields a heat kernel exponent that violates the boundedness requirement from Corollary \ref{cor:intbetatruncated}. This singularity imposes the introduction of fractional derivatives.

Fix $1\leq p<\infty.$ We apply $\Lambda^\frac{1}{2}$ to \eqref{eq:NSDuhamel} and take $L^p$ norm. Proceeding as before, the linear term can be bounded using Remark \ref{rem:heatsmallXphigherorder}. For the nonlinear term, we can repeat the previous strategy as long as we choose the different norms properly. More specifically, we can take
$$
\begin{aligned}
\|\Lambda^\frac{1}{2}\w(t)\|_{L^p}\leq& \|\Lambda^\frac{1}{2}\heat{(t-\sigma)}\w(\sigma)\|_{L^p}+\int_\sigma^t\|\na\Lambda^{\frac{1}{2}}\heat{(t-\tau)}(\w\otimes u-u\otimes \w)(\tau)\|_{L^p}d\tau\\
\lesssim & t^{\frac{1}{p}-\frac{5}{4}}\|\omega\|_{X_t^{\frac{4}{3}}}+\int_\sigma^t (t-\tau)^{\frac{1}{p}-\frac{1}{p}-\frac{3}{4}}\|\w(\tau)\|_{L^{\frac{p(p+2)}{2}}}\|u(\tau)\|_{L^{p+2}}d\tau\\
\lesssim & t^{\frac{1}{p}-\frac{5}{4}}\|\omega\|_{X_t^{\frac{4}{3}}}+ t^{\frac{1}{p}-\frac{5}{4}}\|\omega\|_{X_t^{\frac{p(p+2)}{2}}}\|\omega\|_{X_t^{\frac{2p+4}{p+4}}}\int_{\frac{1}{3}}^1 (1-z)^{-\frac{3}{4}}z^{\frac{1}{p}-\frac{3}{2}}dz,
\end{aligned}$$
to obtain the desired result for $s=1/2$.

We now assume that the result is proved for any $1\leq p<\infty$ and $s=k/2$, for some $k\in \N$. We proceed by induction in $k$.  We apply $\Lambda^{\frac{k+1}{2}}$ to \eqref{eq:NSDuhamel}. To deal with the nonlinear term, we apply $k/2$ derivatives to $\omega\otimes u-u\otimes \w$, and $1/2$ to the heat kernel. We note that, as in Remark \ref{uregularity}, for any $s\geq 0$ and $1<p<2$, with $q$ defined by
$
\frac{1}{q}=\frac{1}{p}-\frac{1}{2},
$
we have
\begin{equation}\label{eq:bounduwgeneral}
\|\Lambda^s u\|_{L^q}\lesssim_q\|\Lambda^s\w\|_{L^p}.
\end{equation} Then, using Lemma \ref{lem:KatoPonce} we can conclude 

$$
\begin{aligned}
\|\Lambda^\frac{k+1}{2}\w(t)\|_{L^p}\leq& \|\Lambda^\frac{k+1}{2}\heat{(t-\sigma)}\w(\sigma)\|_{L^p}+\int_\sigma^t\|\na\Lambda^{\frac{1}{2}}\heat{(t-\tau)}\Lambda^\frac{k}{2}(\w\otimes u-u\otimes \w)(\tau)\|_{L^p}d\tau\\
\lesssim & t^{\frac{1}{p}-1-\frac{k+1}{4}}\|\omega\|_{X_t^{\frac{4}{3}}}+\int_\sigma^t (t-\tau)^{\frac{1}{p}-\frac{1}{p}-\frac{3}{4}}\big(\|\Lambda^\frac{k}{2}\w(\tau)\|_{L^{\frac{p(p+2)}{2}}}\|u(\tau)\|_{L^{p+2}}\\
&\hspace{5cm}+\|\w(\tau)\|_{L^{\frac{p(p+2)}{2}}}\|\Lambda^\frac{k}{2}u(\tau)\|_{L^{p+2}}\big)d\tau\\
\lesssim & t^{\frac{1}{p}-1-\frac{k+1}{4}}\|\omega\|_{X_t^{\frac{4}{3}}}+ t^{\frac{1}{p}-1-\frac{k+1}{4}}\int_{\frac{1}{3}}^1 (1-z)^{-\frac{3}{4}}z^{\frac{1}{p}-\frac{3}{2}-\frac{k}{4}}dz.
\end{aligned}$$
By interpolation, we obtain the result for any $s\in(0,\infty)$. 
For $p=\infty$, we apply Lemma \ref{lem:gagliardoniremberg} to
$f=\Lambda^s\omega$. Using also the boundedness of the Riesz
transforms in $L^4$, we obtain
$$
\begin{aligned}
t^{1+\frac s2}\|\Lambda^s\omega(t)\|_{L^\infty}
&\lesssim
\left(
t^{\frac34+\frac s2}
\|\Lambda^s\omega(t)\|_{L^4}
\right)^{\frac12}
\left(
t^{\frac54+\frac s2}
\|\Lambda^{s+1}\omega(t)\|_{L^4}
\right)^{\frac12}\lesssim 1.
\end{aligned}
$$
\end{proof}

\end{document}
